% ============================================================
%  MỞ ĐẦU
%  Lưu ý: Đây KHÔNG phải chương đánh số – dùng \chapter*
% ============================================================

\setlength{\parskip}{10pt}

\noindent\textbf{1. Đặt vấn đề}

[Trình bày bối cảnh thực tiễn và lý luận dẫn đến sự cần thiết của đề tài.
Nêu rõ vấn đề còn tồn tại mà luận văn hướng đến giải quyết.
Tham khảo các nghiên cứu liên quan để chỉ ra khoảng trống nghiên cứu (research gap).]

\medskip
\noindent\textbf{2. Mục tiêu nghiên cứu}

[Nêu mục tiêu tổng quát và các mục tiêu cụ thể của luận văn.
Mỗi mục tiêu cần đo lường được (SMART).

Ví dụ:
\begin{itemize}
  \item Xây dựng mô hình \ldots\ đạt độ chính xác $\geq$ XX\% trên tập dữ liệu \ldots
  \item Đề xuất phương pháp \ldots\ cải thiện \ldots
  \item Đánh giá hiệu quả \ldots\ so với các phương pháp hiện có.
\end{itemize}]

\medskip
\noindent\textbf{3. Đối tượng và phạm vi nghiên cứu}

\textit{Đối tượng nghiên cứu:} [Xác định cái gì được nghiên cứu – mô hình, thuật toán, hệ thống, dữ liệu...]

\textit{Phạm vi nghiên cứu:}
\begin{itemize}
  \item \textbf{Phạm vi nội dung:} [Chủ đề, bài toán, phương pháp cụ thể]
  \item \textbf{Phạm vi không gian:} [Tập dữ liệu, địa bàn áp dụng]
  \item \textbf{Phạm vi thời gian:} [Giai đoạn dữ liệu; thời gian thực hiện đề tài]
\end{itemize}

\medskip
\noindent\textbf{4. Cách tiếp cận và phương pháp nghiên cứu}

[Mô tả phương pháp luận tổng quát: thực nghiệm, định lượng, học máy...
Nêu ngắn gọn quy trình: thu thập dữ liệu → tiền xử lý → xây dựng mô hình → đánh giá.]

\medskip
\noindent\textbf{5. Ý nghĩa thực tiễn của đề tài}

[Nêu đóng góp của đề tài về mặt khoa học và thực tiễn.
Đặc biệt nhấn mạnh ứng dụng vào bối cảnh Việt Nam hoặc lĩnh vực cụ thể.]
